\section{The ConvergenceSweep Bridge}
\label{sec:bridge}

\subsection{Two Notions of Convergence}

The term ``convergence'' is used differently in the ensemble and algorithmic
settings.  For ensemble methods, convergence refers to the Markov chain
approaching its stationary distribution.  For \cs{} \citep{dellaLibera2026b16},
convergence refers to the optimisation search reaching the minimum edge-cut
map for a given state and seat count.

These are distinct phenomena, but they interact.  The key question is:

\begin{keybox}
At the threshold $T = 600$, \cs{} produces the minimum-edge-cut (\ar{}) plan.
This plan is a \emph{specific point} in the plan space that the \recom{} chain
also explores.  Where does this specific point fall in the \recom{} distribution?
\end{keybox}

\subsection{Characterising the Local Optimum}

The \ar{} plan is a local optimum of the METIS edge-cut objective with
respect to small perturbations.  Under \recom{} moves, the plan will not
remain in its basin of attraction for long---a single recombination move
typically changes $5$--$15\%$ of district assignments and can escape local
optima.  This means the \ar{} plan is \emph{not} the typical plan in the
\recom{} ensemble.

However, the \ar{} plan's \emph{compactness} (as measured by Polsby-Popper
score) is empirically high: T.4 shows a mean $\bar\pp \approx 0.38$--$0.45$
across states, versus \recom{} ensemble means of $0.28$--$0.35$.  This
places the \ar{} plan typically at the $65$th--$75$th percentile of the
\recom{} compactness distribution---more compact than average but not extreme.

\subsection{The Minimum-EC Plan in the ReCom Distribution}

Let $P^*_{\EC}$ denote the minimum-edge-cut plan found by \cs{}.  The
edge-cut objective $\EC(P)$ is negatively correlated with Polsby-Popper
compactness (lower edge cut implies fewer shared-boundary census tracts
across district lines, which correlates with more cohesive district shapes).
Within the \recom{} ensemble, plans with lower $\EC$ are rarer, because
\recom{} does not target $\EC$ minimisation.

Empirically, we estimate $P^*_{\EC}$ falls at approximately the
\begin{itemize}
  \item $65$th--$75$th percentile of the $\bar\pp$ distribution,
  \item $45$th--$60$th percentile of the $S_D$ (partisan seats) distribution,
  \item $50$th--$65$th percentile of the $\mathrm{MM}$ distribution,
\end{itemize}
with state-specific variation characterised in G.1--G.3.  The partisan
percentile near the median is the central finding: the \ar{} plan's
partisanship is close to what geography determines, not to what a partisan
mapmaker would achieve.

\subsection{Diagnostic Consistency}

The $T = 600$ threshold for \cs{} was certified in U.1 using empirical data
from B.7 (50-state seed sweep).%
\footnote{The $T=600$ ConvergenceSweep formula, its theoretical foundations,
and its empirical margin are documented in the companion U.1 paper.}
The same data supports ensemble diagnostic
thresholds: if the seed sweep finds no improvement after 600 seeds, the
plan is stable under seed perturbation.  This is analogous to---but distinct
from---an ensemble chain achieving \ess{} $> 500$ (recommended; minimum floor is 400).

ConvergenceSweep $T = 600$ and ensemble \ess{} are independent certificates
of robustness: \ess{} measures MCMC mixing in the ensemble space; $T = 600$
measures monotonic convergence in the seed-quality space.  Both certify
stability but via orthogonal mechanisms.  They are not numerically comparable
and should not be conflated.  Together they provide complete coverage of the
two-method framework: \ess{} certifies that the ensemble distribution estimate
is stable under additional draws; $T = 600$ certifies that the deterministic
search has found the best available plan under additional seeds.

\subsection{Governed Block-Level Kernel Sensitivity}

The separately frozen Rhode Island Stage~1 gate provides a concrete warning
against treating ``the ReCom ensemble'' as a kernel-free object.  On the same
25,649-block RCTX, starting assignment, $0.005$ population tolerance, chain
schedule, and seed derivation, four Wilson-tree chains and four
random-weight-Kruskal-tree chains each completed 2,000 steps.  After the frozen
500-step burn-in, both cut metrics passed split-$\rhat<1.05$ and pooled
$\ess\geq100$ for each kernel.  A fresh sequential execution regenerated all
normalized metrics and canonical snapshots exactly.

Those execution checks did not make the distributions interchangeable.  Mean
unweighted cut fraction was $0.0103331$ for Wilson and $0.00778270$ for
Kruskal, with descriptive two-sample KS statistic $0.405833$.  Mean
RCTX-weighted boundary cut was $193{,}294{,}034$ and $152{,}599{,}928$,
respectively, with KS statistic $0.31$.  Because the draws are autocorrelated,
the KS $p$-values are descriptive only.  The result supports a Rhode Island
kernel-sensitivity statement; it does not prove mixing, sampler equivalence,
or national representativeness.  A subsequent exact resource audit passed and
mechanically authorized a 21-hour runner budget, a 2.25~GiB per-process memory
ceiling, and 3~GiB retained and scratch ceilings.  Under the separately frozen
NH/NM/GA expansion, five of six primary State--kernel runs completed before the
GA Kruskal process encountered host disk exhaustion.  The registered stop rule
closed that protocol without retry, and no governed replay ran.  The retained
partial package is therefore failure-envelope evidence, not a completed
multi-state comparison.  Any future expansion requires a new protocol with
actual-volume capacity admission at the process-creation boundary; it cannot
continue the failed schedule.
